\documentclass[11pt]{article}

\usepackage[a4paper,margin=1in]{geometry}
\usepackage{amsmath,amssymb,mathtools}
\usepackage[hidelinks]{hyperref}

\title{Single-Set Quantitative Selection Principle}
\author{}
\date{}

\begin{document}
\maketitle

This note implements the concrete selection-principle step from
\texttt{plan1.md}. The purpose is not to replace the whole
Brasco--De Philippis--Velichkov proof, but to isolate the part that can be
made quantitative without changing their free-boundary machinery.

The main output is a single-set form of BDV Proposition 4.4. Starting from one
near-counterexample, it constructs a selected quasi-open minimizer whose
energy-to-asymmetry ratio is controlled up to constants.

\section{Normalization}

Use the BDV torsion convention
\[
E(\Omega)
=
\min_{u\in H^1_0(\Omega)}
\left\{
\frac12\int_\Omega |\nabla u|^2-\int_\Omega u
\right\}
=
-\frac12\int_\Omega u_\Omega.
\]
Fix
\[
|\Omega|=|B_1|=\omega_N.
\]
Then the Saint-Venant deficit is
\[
\delta(\Omega):=E(\Omega)-E(B_1)\ge0.
\]

For the selection step, use BDV's smoother asymmetry
\[
\alpha(\Omega)
:=
\int_{\Omega\Delta B_1(x_\Omega)}
\bigl|1-|x-x_\Omega|\bigr|\,dx,
\]
where \(x_\Omega\) is the barycenter. BDV Lemma 4.2 gives, for bounded
\(\Omega\subset B_R\),
\[
|\Omega\Delta B_1(x_\Omega)|^2\lesssim_N \alpha(\Omega),
\]
\[
|\alpha(\Omega_1)-\alpha(\Omega_2)|
\lesssim_R
|\Omega_1\Delta\Omega_2|,
\]
and, for nearly spherical sets,
\[
\alpha(\Omega)\lesssim_N
\|\varphi\|_{L^2(\partial B_1)}^2.
\]
Thus the right badness ratio for the selection step is
\[
Q_\alpha(\Omega)
:=
\frac{E(\Omega)-E(B_1)}{\alpha(\Omega)}.
\]
The final Fraenkel statement is recovered later from
\(\mathcal A(\Omega)^2\lesssim \alpha(\Omega)\).

\section{Penalized Problem}

Let \(\Omega_0\subset B_R\), \(R\ge2\), satisfy
\[
|\Omega_0|=|B_1|,
\qquad
\varepsilon:=\alpha(\Omega_0)>0,
\qquad
\delta:=E(\Omega_0)-E(B_1),
\qquad
q:=\delta/\varepsilon.
\]
Assume \(q\ll1\). Choose
\[
\tau:=q^{1/4},
\]
or more generally any \(\tau\) with
\[
q\ll \tau^2\ll1.
\]
Let \(f_{\hat\eta}\) be the BDV volume penalty from Lemma 4.5:
\[
f_{\hat\eta}(s)=
\begin{cases}
\hat\eta(s-\omega_N),&s\le\omega_N,\\
(s-\omega_N)/\hat\eta,&s\ge\omega_N.
\end{cases}
\]
Set
\[
F_{\hat\eta}(U):=E(U)+f_{\hat\eta}(|U|).
\]
BDV Lemma 4.5 says that, for \(\hat\eta=\hat\eta(N,R)\), \(B_1\) is a
minimizer of \(F_{\hat\eta}\) among sets in \(B_R\).

Now minimize over quasi-open \(U\subset B_R\):
\[
\mathcal G_{\varepsilon,\tau}(U)
:=
F_{\hat\eta}(U)
+
\sqrt{\varepsilon^2+\tau^2(\alpha(U)-\varepsilon)^2}.
\]
This is exactly BDV's penalized functional (4.10), with the sequence parameter
\(\sigma\) replaced by the single-set parameter \(\tau\).

\section{Quantitative Selected-Minimizer Estimates}

Assume \(\tau\le\tau_0(N,R)\), small enough that the BDV existence and
quasi-minimizer estimates hold with \(\tau\) in place of \(\sigma\). Let
\(U_*\) be a minimizer of \(\mathcal G_{\varepsilon,\tau}\).

\subsection*{Lemma 1. Energy and Asymmetry Preservation Before Normalization}

The minimizer \(U_*\) satisfies
\[
0\le
F_{\hat\eta}(U_*)-F_{\hat\eta}(B_1)
\le
\delta,
\]
and
\[
|\alpha(U_*)-\varepsilon|
\le
\frac{\sqrt{2\varepsilon\delta+\delta^2}}{\tau}.
\]
In particular, if \(q=\delta/\varepsilon\le1\) and \(\tau=q^{1/4}\), then
\[
|\alpha(U_*)-\alpha(\Omega_0)|
\le
C\alpha(\Omega_0)q^{1/4}.
\]

\paragraph{Proof.}
By minimality against \(\Omega_0\),
\[
F_{\hat\eta}(U_*)
+
\sqrt{\varepsilon^2+\tau^2(\alpha(U_*)-\varepsilon)^2}
\le
F_{\hat\eta}(\Omega_0)+\varepsilon.
\]
Since \(|\Omega_0|=|B_1|\),
\[
F_{\hat\eta}(\Omega_0)=E(\Omega_0)
=F_{\hat\eta}(B_1)+\delta.
\]
Because \(B_1\) minimizes \(F_{\hat\eta}\),
\[
F_{\hat\eta}(U_*)\ge F_{\hat\eta}(B_1).
\]
Therefore
\[
F_{\hat\eta}(U_*)-F_{\hat\eta}(B_1)\le\delta,
\]
and
\[
\sqrt{\varepsilon^2+\tau^2(\alpha(U_*)-\varepsilon)^2}
\le
\varepsilon+\delta.
\]
Squaring the last inequality gives
\[
\tau^2(\alpha(U_*)-\varepsilon)^2
\le
2\varepsilon\delta+\delta^2.
\]
This proves the claim.

\subsection*{Lemma 2. Volume Error}

There is \(C=C(N,R)\) such that
\[
\bigl||U_*|-|B_1|\bigr|
\le
C\delta.
\]

\paragraph{Proof.}
Let \(B_{r_*}\) be a ball with \(|B_{r_*}|=|U_*|\). By rearrangement,
\[
F_{\hat\eta}(B_{r_*})\le F_{\hat\eta}(U_*).
\]
BDV Lemma 4.5 gives
\[
F_{\hat\eta}(B_{r_*})-F_{\hat\eta}(B_1)
\ge
\frac{|r_*-1|}{C_4(N,R)}.
\]
Together with Lemma 1,
\[
|r_*-1|\le C\delta.
\]
Since \(|B_{r_*}|=\omega_N r_*^N\), this implies
\[
\bigl||U_*|-|B_1|\bigr|\le C\delta.
\]

\subsection*{Lemma 3. Ratio Preservation After Volume Normalization}

Let
\[
\lambda:=\left(\frac{|B_1|}{|U_*|}\right)^{1/N},
\qquad
\widetilde\Omega:=\lambda U_*.
\]
After translating if needed, \(\widetilde\Omega\) has volume \(|B_1|\), and
\[
\alpha(\widetilde\Omega)
\ge
\left(1-Cq^{1/4}\right)\alpha(\Omega_0),
\]
\[
E(\widetilde\Omega)-E(B_1)
\le
C\big(E(\Omega_0)-E(B_1)\big),
\]
provided \(q\le q_0(N,R)\) and \(\alpha(\Omega_0)\) is in the small regime.
Consequently,
\[
Q_\alpha(\widetilde\Omega)
\le
C Q_\alpha(\Omega_0).
\]

\paragraph{Proof sketch.}
Lemma 2 gives
\[
|\lambda-1|\le C\delta=C\varepsilon q.
\]
Since \(q\ll1\), this error is smaller than the asymmetry-preservation error
\(\varepsilon q^{1/4}\). The selected minimizer has the perimeter bound from
the BDV existence argument, so the small dilation satisfies
\[
|U_*\Delta \lambda U_*|
\le C(N,R)|\lambda-1|.
\]
The Lipschitz property of \(\alpha\) on bounded sets then gives
\[
|\alpha(\widetilde\Omega)-\alpha(U_*)|
\le C\varepsilon q.
\]
Combining with Lemma 1 gives
\[
|\alpha(\widetilde\Omega)-\varepsilon|
\le C\varepsilon q^{1/4}.
\]

For the energy, use
\[
E(\lambda U_*)=\lambda^{N+2}E(U_*),
\]
the bound \(F_{\hat\eta}(U_*)-F_{\hat\eta}(B_1)\le\delta\), the Lipschitz
control of \(f_{\hat\eta}\) in the volume variable, and
\(|\lambda-1|\le C\delta\). Since \(U_*\subset B_R\), the energy is bounded in
terms of \(N,R\), so the dilation changes the deficit by \(O_{N,R}(\delta)\).
Thus
\[
E(\widetilde\Omega)-E(B_1)\le C\delta.
\]
The ratio estimate follows from the two displayed bounds.

\section{Free-Boundary Regularity Inherited from BDV}

The minimizer \(U_*=\{u_*>0\}\) satisfies the perturbed one-phase minimality
inequality
\[
\frac12\int |\nabla u_*|^2-\int u_*
+f_{\hat\eta}(|\{u_*>0\}|)
\le
\frac12\int |\nabla v|^2-\int v
+f_{\hat\eta}(|\{v>0\}|)
+C_R\tau|\{u_*>0\}\Delta\{v>0\}|
\]
for all \(v\in H^1_0(B_R)\). This is BDV (4.19) with \(\tau\) in place of
\(\sigma\).

Hence, for \(\tau\le\tau_0(N,R)\), the Alt-Caffarelli estimates used in BDV
apply:
\begin{itemize}
\item nondegeneracy of \(u_*\) near the free boundary,
\item Lipschitz control,
\item density estimates for both phases,
\item reduced-boundary structure,
\item smoothness once the free boundary is sufficiently flat.
\end{itemize}

The Euler-Lagrange Bernoulli coefficient has the same structure as BDV Lemma
4.16:
\[
q_{u_*}(x)^2
=
\Lambda
+
\frac{\tau^2(\alpha(U_*)-\varepsilon)}
{\sqrt{\varepsilon^2+\tau^2(\alpha(U_*)-\varepsilon)^2}}
\cdot
\Psi_{U_*}(x),
\]
where \(\Psi_{U_*}\) is the smooth \(\alpha\)-variation term. The prefactor is
bounded by \(\tau\), so \(q_{u_*}\) has uniform \(C^k\) bounds near
\(\partial U_*\), exactly as in BDV Lemma 4.16.

\section{Quantitative Entry into Hausdorff Closeness}

The compactness step in BDV can be partially quantified as follows.

\subsection*{Lemma 4. From \(\alpha\) and Density to Hausdorff Distance}

Assume \(U\subset B_R\) satisfies the two-sided density estimates
\[
|U\cap B_r(x)|\ge c r^N,\qquad
|B_r(x)\setminus U|\ge c r^N
\]
for every \(x\in\partial U\) and every \(r\le r_0\). If, after centering,
\[
|U\Delta B_1|\le M,
\]
then
\[
d_H(\partial U,\partial B_1)
\le
C(N,R,c,r_0) M^{1/N}.
\]
Using BDV Lemma 4.2,
\[
M\lesssim_N \alpha(U)^{1/2},
\]
so selected minimizers satisfy
\[
d_H(\partial U,\partial B_1)
\le
C\alpha(U)^{1/(2N)}.
\]

\paragraph{Proof sketch.}
If \(x\in\partial U\) and \({\rm dist}(x,\partial B_1)=d\), then for
\(r\sim d\) the ball \(B_r(x)\) lies either inside \(B_1\) or outside \(B_1\).
The density estimate forces a chunk of size \(\gtrsim r^N\) in
\(U\Delta B_1\). Thus \(d^N\lesssim M\).

Conversely, if \(y\in\partial B_1\) is distance \(d\) from \(\partial U\), then
\(B_d(y)\) is either contained in \(U\) or disjoint from \(U\). Either way, a
fixed fraction of \(B_d(y)\) lies in \(U\Delta B_1\), so again
\(d^N\lesssim M\).

\section{Remaining Quantitative Gap}

BDV's final regularity passage is
\[
\begin{aligned}
L^1\hbox{ closeness}
&\Rightarrow
\hbox{Kuratowski/Hausdorff closeness} \\
&\Rightarrow
\hbox{flatness at a fixed scale} \\
&\Rightarrow
C^{1,\gamma}\hbox{ free-boundary graph}
\Rightarrow
C^\infty\hbox{ convergence}.
\end{aligned}
\]
Lemma 4 above supplies the Hausdorff step quantitatively for selected
minimizers. The remaining task is to make the Alt-Caffarelli flatness threshold
explicit enough to state
\[
\alpha(U_*)\le \varepsilon_0(N,R,\tau)
\quad\Longrightarrow\quad
\partial U_*
\hbox{ is a global }C^{1,\gamma}\hbox{ graph over }\partial B_1.
\]
Once this graph entry is available, BDV's nearly spherical second variation
gives
\[
E(\widetilde\Omega)-E(B_1)
\gtrsim_N
\|\varphi\|_{H^{1/2}(\partial B_1)}^2
\gtrsim_N
\alpha(\widetilde\Omega),
\]
contradicting \(Q_\alpha(\widetilde\Omega)\ll1\).

Thus the selection route has been reduced to one concrete missing quantitative
regularity statement: track the flatness constants in the Alt-Caffarelli
improvement step for the selected minimizers.

\end{document}
